\begin{problem}{TST 2026 - Bài 3}{standard input}{standard output}{2 seconds}{256 megabytes}

Hôm nay, Alice và Bob sẽ thực hiện một trò ảo thuật trước mặt hàng trăm khán giả. Trò ảo thuật mà hai người sẽ làm có thể mô tả như sau:
\begin{itemize}
   \item Đầu tiên, Bob bị bịt mắt, Alice gọi một khán giả bất kỳ lên. Khán giả sẽ chọn $N$ là bài trong bộ bài Tây $52$ lá Alice đã chuẩn bị.
   \item Từ $N$ lá bài, khán giả vẽ lên bảng một cây \textbf{có gốc} gồm $N$ đỉnh, mỗi đỉnh tương ứng với một lá bài.
   \item Tiếp theo, Alice sẽ viết một số nguyên dương ở sau tất cả lá bài khán giả đã chọn.
   \item Sau đó, Alice xáo lại các lá bài và đưa cho Bob đống bài vừa xáo. Bob không được nhìn mặt trước của các lá bài mà chỉ được nhìn các số nguyên dương mà Alice đã viết trên đó.
   \item Từ dữ kiện trên, Bob sẽ dựng lại cây mà khán giả đã vẽ. Trò ảo thuật sẽ được coi là thành công nếu như cây mà Bob vẽ giống hệt với cây mà khán giả đã vẽ ban đầu. Cụ thể:
   \begin{itemize}
      \item Hai lá bài được chọn làm gốc phải giống hệt nhau.
      \item Các cặp lá bài có cạnh nối trực tiếp mô phỏng quan hệ cha-con trên cây cũng phải giống hệt nhau. Nếu trên cây của khán giả, lá bài $A$ là cha của lá bài $B$ nhưng trên cây của Bob, lá bài $B$ là cha của lá bài $A$ thì hai cây sẽ không được coi là giống nhau.
      \item \textbf{Lưu ý:} Hai lá bài được cho là giống nhau nếu mặt trước của hai lá bài giống hệt nhau.
   \end{itemize}
\end{itemize}

Vì giới hạn kích thước của bảng, cây mà khán giả vẽ luôn đảm bảo có bậc tối đa nhỏ hơn hoặc bằng $8$ và số đỉnh tối đa có cùng bậc nhỏ hơn hoặc bằng $8$ $^\dagger$. Trong trò ảo thuật này, cây sẽ được mô phỏng dưới dạng một danh sách kề, mỗi đỉnh lưu trữ tất cả các con kề với đỉnh đó. Gốc của cây là đỉnh duy nhất không nằm trong danh sách kề của bất kỳ đỉnh nào.

Trò ảo thuật sẽ khiến mọi người càng trầm trồ và thán phục nếu giá trị lớn nhất mà Alice viết trên các lá bài càng nhỏ.

\textbf{Yêu cầu:} Hãy viết một chương trình mô phỏng lại trò ảo thuật của Alice và Bob.

$\dagger$: Trong cây, một đỉnh có bậc là $k$ nếu trên đường đi ngắn nhất từ gốc đến đỉnh đó có đúng $k$ cạnh. Bậc tối đa của cây được coi là bậc tối đa của các đỉnh trên cây đó.

\InputFile
\textbf{Chi tiết cài đặt}

Thí sinh cần cài đặt hai hàm sau:

\begin{lstlisting}
vector<int> solveAlice(vector<vector<int>> adj);
\end{lstlisting}

\begin{itemize}
   \item $adj$: Một mảng hai chiều mô phỏng danh sách kề của cây mà khán giả đã vẽ. Trong đó:
   \begin{itemize}
      \item $N=|adj|$
      \item Các đỉnh trên cây được đánh số từ $0$ đến $N-1$, trong đó đỉnh thứ $i$ ($0\le i\le N-1$) tương ứng với lá bài thứ $i$
      \item Mỗi phần tử $adj_i$ ($0\le i\le N-1$) là một mảng lưu tất cả các \textbf{đỉnh con} kề với đỉnh $i$ trên cây mà khán giả đã vẽ
      \item Giá trị duy nhất không xuất hiện trên các mảng $adj_i$ là gốc của cây.
      \item $\sum_{|adj_i|}=N-1$
   \end{itemize}
   \item Hàm này cần trả về một mảng $S$ gồm $N$ phần tử, trong đó $S_i$ ($0\le i\le N-1$) là số nguyên dương Alice đánh trên lá bài thứ $i$. Các số nguyên dương Alice viết không nhất thiết phải đôi một phân biệt.
\end{itemize}

\begin{lstlisting}
vector<vector<int>> solveBob(vector<int> S);
\end{lstlisting}

\begin{itemize}
   \item $S$: Một mảng gồm $N$ phần tử, trong đó $S_i$ ($0\le i\le N-1$) đại diện cho giá trị được viết trên lá bài thứ $i$ từ đống bài Alice đưa cho Bob. Các giá trị này có thể được sắp xếp lại theo một thứ tự bất kỳ, không nhất thiết phải theo thứ tự các giá trị xuất hiện trên mảng mà Alice đã trả về.
   \item Bob sẽ chỉ được biết các số được đánh trên các lá bài mà không được biết mặt trước các lá bài Alice đưa trông như thế nào.
   \item Hàm này cần trả về một mảng hai chiều tương ứng với danh sách kề của cây mà Bob vẽ. Danh sách kề này cần mô phỏng cây y hệt so với cây mà khán giả đã vẽ, với các chỉ số trong danh sách tương ứng với trình tự các lá bài mà Alice đã xáo.
\end{itemize}

Trình chấm của bài toán này sẽ hoạt động như sau:
\begin{itemize}
   \item Trong một test, trình chấm sẽ chạy $T$ lần, mỗi lần chạy tương ứng với một trường hợp thử nghiệm. Trong một lần chạy:
   \begin{itemize}
      \item Đầu tiên, trình chấm gọi hàm \texttt{solveAlice()} trước. Thông tin của mảng $S$ từ hàm \texttt{solveAlice()} sẽ được trình chấm lưu trữ lại.
      \item Sau đó, trình chấm sẽ gọi hàm \texttt{solveBob()}. Khi đã nhận về kết quả từ hàm \texttt{solveBob()}, trình chấm sẽ kiểm tra tính hợp lệ của kết quả. Nếu đáp án đưa ra thỏa mãn điều kiện đề bài, trình chấm sẽ thực hiện lần chạy tiếp theo, hoặc dừng lại khi đã thực hiện xong tất cả $T$ lần chạy.
   \end{itemize}
   \item \textbf{Lưu ý: Các lần gọi hàm \texttt{solveAlice()} và \texttt{solveBob()} được gọi trên các chương trình hoàn toàn riêng biệt}.
   \item Sau khi chương trình đã thực hiện đủ $T$ lần chạy, trình chấm sẽ tính điểm dựa trên độ tối ưu chương trình của bạn (sẽ nói rõ ở phần cách tính điểm).
\end{itemize}

\textbf{Lưu ý:}
\begin{itemize}
   \item Chương trình của bạn cần phải khai báo thư viện \texttt{"magicianlib.h"} ở đầu.
   \item Trong chương trình, thí sinh được phép cài đặt các biến, các hàm toàn cục, nhưng không được phép có hàm \texttt{main()}.
   \item Chương trình của bạn không được tương tác trực tiếp với đầu vào chuẩn (stdin) và đầu ra chuẩn (stdout).
\end{itemize}

\textbf{Ràng buộc}
\begin{itemize}
   \item $T\le 100$
   \item $N\le 32$
   \item Cấu trúc của cây luôn đảm bảo bậc tối đa của cây nhỏ hơn hoặc bằng $8$ và số đỉnh tối đa có cùng bậc nhỏ hơn hoặc bằng $8$
\end{itemize}

\Scoring
\begin{center}
  \begin{tabular}{|c|c|l|} \hline
    \textbf{Subtask} & \textbf{Điểm} & \textbf{Giới hạn} \\ \hline
    1 & $10$ & $N\le 9$, cây mà khán giả vẽ luôn đảm bảo là một đường thẳng \\ \hline
    2 & $90$ & Không có ràng buộc gì thêm \\ \hline
  \end{tabular}
\end{center}

\textbf{Cách tính điểm}

Nếu trong bất kỳ trường hợp thử nghiệm nào, hàm \texttt{solveBob()} trả về danh sách kề không thỏa mãn với điều kiện đề bài, bạn sẽ không nhận được điểm cho test đó.

Ngược lại, gọi $M$ là số nguyên dương lớn nhất được viết trên các lá bài của Alice trong tất cả trường hợp thử nghiệm. Khi đó điểm của bạn được tính theo công thức như sau:

\begin{center}
  \begin{tabular}{|c|c|} \hline
    \textbf{Giá trị $M$} & \textbf{Phần trăm số điểm} \\ \hline
    $M\le 56$ & $100\%$ \\ \hline
    $M=57$ & $80\%$ \\ \hline
    $58\le M\lt 88$ & $(176-2M)\%$ \\ \hline
    $M\ge 88$ & $0\%$ \\ \hline
  \end{tabular}
\end{center}

Điểm của một subtask bằng phần trăm điểm tối thiểu của một test trong subtask đó nhân với số điểm tối đa của subtask.


\end{problem}

